% ==========================================================================
%  La hora dorada -- informe Fase 5
%  Tipografia: Arial (Helvetica, sustituto estandar de Arial en pdfLaTeX),
%  cuerpo 12 pt, texto justificado.
% ==========================================================================
\documentclass[12pt,a4paper]{article}

\usepackage[main=spanish,provide=*]{babel}
\usepackage[utf8]{inputenc}
\usepackage[T1]{fontenc}
\usepackage[scaled=0.95]{helvet}           % Arial (Helvetica: sustituto estandar en pdfLaTeX)
\renewcommand{\familydefault}{\sfdefault}
\usepackage[margin=2cm]{geometry}
\usepackage{graphicx}
\usepackage{array}                         % columnas p{} con \raggedright
\usepackage{booktabs}
\usepackage{float}                         % [H]: cada cuadro donde se lo escribe, sin solaparse
\usepackage{xcolor}
\usepackage{url}
\usepackage[hidelinks]{hyperref}

% Rutas y URLs largas: que puedan cortarse de linea en cualquier caracter
\urlstyle{tt}
\def\UrlBreaks{\do\/\do\-\do\.\do\_\do\:\do\?\do\=\do\&\do\#%
  \do\a\do\b\do\c\do\d\do\e\do\f\do\g\do\h\do\i\do\j\do\k\do\l\do\m%
  \do\n\do\o\do\p\do\q\do\r\do\s\do\t\do\u\do\v\do\w\do\x\do\y\do\z%
  \do\0\do\1\do\2\do\3\do\4\do\5\do\6\do\7\do\8\do\9}
\Urlmuskip=0mu plus 1mu\relax
\setlength{\emergencystretch}{3em}          % evita cajas desbordadas por palabras largas
% ruta o nombre de archivo: parte SOLO en los separadores, no a mitad de palabra
\newcommand{\rc}[1]{\begingroup
  \def\UrlBreaks{\do\/\do\_\do\-\do\.}\path{#1}\endgroup}
\newcolumntype{R}[1]{>{\raggedright\arraybackslash}p{#1}}

\graphicspath{{figures/}}
\setlength{\parskip}{0.3em plus 0.1em}      % parrafos en bloque (sin sangria)
\setlength{\parindent}{0pt}
\linespread{1.0}
\newcommand{\FloatBarrier}{\par}            % sin el paquete placeins

% --- Densidad: listas y titulos mas compactos (sin enumitem/titlesec) -------
\makeatletter
\renewcommand\@listi{\leftmargin\leftmargini \topsep 2\p@ \parsep \z@ \itemsep 2\p@}
\let\@listI\@listi
\renewcommand\@listii{\leftmargin\leftmarginii \topsep 1\p@ \parsep \z@ \itemsep 1\p@}
\renewcommand\section{\@startsection{section}{1}{\z@}%
  {-2.0ex \@plus -.5ex \@minus -.2ex}{1.0ex \@plus .2ex}{\normalfont\large\bfseries}}
\renewcommand\subsection{\@startsection{subsection}{2}{\z@}%
  {-1.6ex \@plus -.5ex \@minus -.2ex}{0.5ex \@plus .2ex}{\normalfont\normalsize\bfseries}}
\makeatother

% --- Cuadros: filas separadas por lineas y con algo de aire ----------------
\renewcommand{\arraystretch}{1.15}
\setlength{\tabcolsep}{5pt}
\newcommand{\filasep}{\addlinespace[1pt]\midrule\addlinespace[1pt]}

% --- Insumos generados por src/export.py -------------------------------------
\newcommand{\inputsihay}[1]{\IfFileExists{#1}{\input{#1}}{%
  \begin{center}\itshape[Tabla pendiente: ejecutar \texttt{python src/export.py}]\end{center}}}
\newcommand{\figsihay}[3][0.8\linewidth]{\IfFileExists{figures/#2.pdf}{%
  \begin{figure}[H]\centering\includegraphics[width=#1]{#2.pdf}%
  \caption{#3}\label{fig:#2}\end{figure}}{%
  \begin{center}\itshape[Figura pendiente: ejecutar \texttt{python src/export.py}]\end{center}}}
% dos figuras lado a lado, una sola leyenda
\newcommand{\dosfiguras}[5][0.47\linewidth]{%
  \IfFileExists{figures/#2.pdf}{%
  \begin{figure}[H]\centering
  \includegraphics[width=#1]{#2.pdf}\hfill\includegraphics[width=#1]{#3.pdf}%
  \caption{#4}\label{fig:#2}\end{figure}}{%
  \begin{center}\itshape[Figuras pendientes: ejecutar \texttt{python src/export.py}]\end{center}}}

% --- Cifras: src/export.py escribe tables/_numeros.tex ----------------------
\IfFileExists{tables/_numeros.tex}{\newcommand{\factorRodeo}{1.82}
\newcommand{\umbralOro}{60}
\newcommand{\nCentrosPoblados}{5\,000}
\newcommand{\poblacionTotal}{3\,174\,254}
\newcommand{\pctCubierta}{85.1}
\newcommand{\pctSinRuta}{11.9}
\newcommand{\poblacionSinRuta}{376\,801}
\newcommand{\medianaAcceso}{23}
\newcommand{\giniTotal}{0.708}
\newcommand{\spearmanAltitud}{0.35}
\newcommand{\peorDistrito}{IPARIA}
\newcommand{\fraccionMuestra}{29}
}{}
\providecommand{\umbralOro}{60}
\providecommand{\nCentrosPoblados}{5\,000}
\providecommand{\poblacionTotal}{--}
\providecommand{\pctCubierta}{--}
\providecommand{\pctSinRuta}{--}
\providecommand{\poblacionSinRuta}{--}
\providecommand{\medianaAcceso}{--}
\providecommand{\giniTotal}{--}
\providecommand{\spearmanAltitud}{--}
\providecommand{\peorDistrito}{--}
\providecommand{\factorRodeo}{--}
\providecommand{\fraccionMuestra}{29}
\providecommand{\coordRecuperadas}{--}
\providecommand{\coordFaltantes}{--}
\providecommand{\tasaRecuperacion}{--}

\title{\textbf{La hora dorada}\\[2mm]
\large ?`Cu\'anto tarda la poblaci\'on de Piura, Ayacucho y Ucayali en llegar
por carretera a un hospital que pueda resolver una emergencia?}
\author{Proyecto integrador --- Data Science con Python}
\date{\today}

\begin{document}
\maketitle
\vspace{-1.1cm}

\begin{abstract}
\noindent
Med\'i, para \nCentrosPoblados{} centros poblados de Piura (costa), Ayacucho
(sierra) y Ucayali (Amazon\'ia), cu\'anto se tarda en auto ---por carreteras
reales, no en l\'inea recta--- hasta el establecimiento de salud m\'as
cercano capaz de operar o atender un parto complicado (categor\'ia
\mbox{II-1} o superior, en funcionamiento). De \poblacionTotal{} habitantes
analizados, el \pctCubierta\,\% est\'a a \umbralOro{} minutos o menos de uno
de esos establecimientos. Un \pctSinRuta\,\% (\poblacionSinRuta{} personas)
vive en lugares sin ninguna carretera que los conecte con un hospital
resolutivo; a esa poblaci\'on la dej\'e marcada como ``sin ruta'' y no le
invent\'e una distancia en l\'inea recta. La mediana del tiempo de acceso es
\medianaAcceso{} minutos. La desigualdad del acceso, medida con el
coeficiente de Gini, es \giniTotal{} (0 = todos igual; 1 = m\'axima
desigualdad). El distrito m\'as cr\'itico es \peorDistrito. El hallazgo
central es geogr\'afico: la Amazon\'ia no tiene un problema de tiempo de
viaje, tiene un problema de \emph{medio} de transporte ---se mueve por r\'io,
no por carretera--- y eso lo esconde cualquier promedio que no pondere por
poblaci\'on.
\end{abstract}

\section{Introducci\'on y planteamiento del problema}

En una emergencia grave ---un accidente serio, una hemorragia en el parto---
existe una ventana de tiempo, la ``hora de oro'', en la que el paciente debe
llegar a un centro que pueda \emph{resolver} el problema. En el Per\'u un
puesto de salud (categor\'ia \mbox{I-1}) puede estabilizar a alguien, pero
no operarlo ni hacer una ces\'area. Para eso hace falta un establecimiento de
categor\'ia \mbox{II-1} o superior.

La pregunta que quise responder es sencilla de enunciar y dif\'icil de
contestar bien: \textbf{?`cu\'anto tarda de verdad la gente en llegar por
carretera a uno de esos establecimientos, y d\'onde est\'an las peores
brechas?} La palabra ``de verdad'' es la clave. Una l\'inea recta en el mapa
no es una carretera. Un establecimiento que figura en el registro puede
estar cerrado o mal ubicado. Una coordenada puede venir con el signo
cambiado. El valor del trabajo est\'a, sobre todo, en c\'omo trat\'e esos
problemas, y este informe se dedica a explicarlo.

Eleg\'i tres departamentos con geograf\'ias deliberadamente distintas
---Piura (costa), Ayacucho (sierra) y Ucayali (Amazon\'ia)--- porque el
contraste es justamente lo interesante: no deber\'ian comportarse igual, y el
informe explica por qu\'e. Los tres est\'an declarados en
\texttt{config.md}; cambiarlos ah\'i cambia todo el an\'alisis sin tocar una
sola l\'inea de c\'odigo.

\section{Archivos entregados}

Todo el proyecto son m\'odulos de Python reutilizables; no hay l\'ogica
pegada en cuadernos. La Tabla~\ref{tab:archivos} indica qu\'e archivo hace
cada fase.

\begin{table}[htbp]
\centering
\footnotesize
\caption{Archivos entregados y su rol.}
\label{tab:archivos}
\begin{tabular}{@{}R{2.7cm}R{3.7cm}R{7.6cm}@{}}
\toprule
Fase & Archivo & Qu\'e hace \\
\midrule
Configuraci\'on & \texttt{config.md} & Todos los par\'ametros: departamentos, rutas, categor\'ias resolutivas, umbrales, motor de ruteo. Nada est\'a escrito a mano en el c\'odigo. \\
\filasep
1 -- adquisici\'on & \texttt{src/acquisition.py} & Resuelve cada fuente (usa la copia local, la copia de \rc{_data/} o la descarga) y controla el encoding al leer. Registra todo en \rc{data/raw/download_log.json}. \\
\filasep
1 -- validaci\'on & \texttt{src/validation.py} & Limpieza, reglas de calidad, recorte a los tres departamentos y guardado en GeoPackage. Escribe \rc{logs/quality_report.csv}. \\
\filasep
1 -- mapas & \texttt{src/mapas.py} & Un mapa interactivo por departamento (\rc{data/outputs/mapa_acceso_*.html}). \\
\filasep
2 -- ruteo & \texttt{src/routing.py} & Construye el grafo vial desde el \texttt{.pbf} local, engancha los puntos, calcula la matriz de tiempos y la cachea en Parquet. \\
\filasep
3 -- poblaci\'on & \texttt{src/poblacion.py} & Estima la poblaci\'on de cada centro poblado a partir de WorldPop. \\
\filasep
3 -- m\'etricas & \texttt{src/metrics.py} & Cada indicador es una funci\'on que recibe una tabla y devuelve una tabla. Exporta a \rc{data/outputs/*.csv}. \\
\filasep
4 -- tablero & \texttt{app.py} & Tablero Streamlit; lee solo archivos ya calculados. \\
\filasep
5 -- informe & \texttt{src/export.py}, \texttt{report/main.tex} & Genera las figuras y tablas del informe y este documento. \\
\bottomrule
\end{tabular}
\end{table}

\FloatBarrier
\section{Fuentes de datos}

\begin{table}[htbp]
\centering
\footnotesize
\caption{Fuentes utilizadas, con fecha de acceso, licencia y limitaciones conocidas.}
\label{tab:fuentes}
\begin{tabular}{@{}R{2.5cm}R{2.9cm}R{2.7cm}R{5.9cm}@{}}
\toprule
Fuente & Contenido & Acceso / licencia & Limitaciones conocidas \\
\midrule
RENIPRESS / SUSALUD & Registro nacional de establecimientos (la oferta) &
CSV mensual, versi\'on del 2026-04-30; datos abiertos &
Categor\'ias sin estandarizar; $\sim$37\,\% de registros sin coordenada;
``activo'' no garantiza personal ni insumos. El registro cambia mes a mes
(el establecimiento 00003598 de Huamanga era \mbox{I-4} en abril y
\mbox{II-1} en agosto de 2026). \\
\filasep
SIGMED / MINEDU & Centros poblados con coordenadas (la demanda) &
Descarga espacial (shapefile); uso educativo &
No trae poblaci\'on; incluye centros poblados sin ninguna v\'ia de acceso. \\
\filasep
OpenStreetMap (extracto de Per\'u, Geofabrik) & Red vial &
\rc{peru-latest.osm.pbf}, 2026-09-07; licencia ODbL &
Cobertura despareja: buena en ciudades, escasa en zona rural y amaz\'onica. \\
\filasep
INEI --- cartograf\'ia censal & L\'imites de distrito, provincia y departamento &
GeoPackage; uso oficial &
La precisi\'on del pol\'igono var\'ia en zonas remotas. \\
\filasep
WorldPop 2020 (1\,km, ajustado a totales ONU) & Poblaci\'on por celda &
GeoTIFF; licencia CC-BY 4.0; \rc{DOI:10.5258/SOTON/WP00685} &
Resoluci\'on de 1\,km; es un modelo, no un censo. \\
\bottomrule
\end{tabular}
\end{table}

\paragraph{Por qu\'e WorldPop y no la poblaci\'on del censo por centro
poblado.} La capa de centros poblados (SIGMED) no incluye poblaci\'on, y solo
el 64\,\% de sus puntos trae el c\'odigo del INEI para poder cruzarlos con la
tabla censal; el 36\,\% restante son puntos agregados por el IGN, sin
enumeraci\'on censal. WorldPop resuelve las dos cosas: cubre el 100\,\% de
los puntos y est\'a ajustado a los totales oficiales, de modo que al sumar
por departamento coincide con el censo (Piura $\approx$ 1{,}86 millones;
Ayacucho $\approx$ 0{,}74; Ucayali $\approx$ 0{,}57). Es, adem\'as, un
producto pensado justamente para estimar denominadores de poblaci\'on en
estudios de acceso a salud.

\FloatBarrier
\section{Metodolog\'ia}

\subsection{Qu\'e cuenta como ``resolutivo''}
Consider\'e resolutivo un establecimiento solo si cumple las dos condiciones:
(i)~su estado es \texttt{ACTIVO} y (ii)~su categor\'ia, ya normalizada, es
\mbox{II-1}, \mbox{II-2}, \mbox{II-E}, \mbox{III-1}, \mbox{III-2} o
\mbox{III-E}. Los de categor\'ia \mbox{I-1} a \mbox{I-4} se quedan en la base
de datos pero no entran en el c\'alculo del m\'as cercano. La categor\'ia en
RENIPRESS viene escrita de muchas maneras (``II 1'', ``II-1'', ``II1''), as\'i
que la reescribo a una forma \'unica con una regla a la vista en el c\'odigo
(\texttt{normalizar\_categoria} en \texttt{src/validation.py}); lo vac\'io,
``0'' o ``S/C'' queda como categor\'ia desconocida.

\subsection{Limpieza y validaci\'on}
La idea de fondo: \textbf{ninguna fila se descarta en silencio}. Cada regla
deja registro de cu\'antos casos afect\'o, qu\'e se hizo y por qu\'e
(Tabla~\ref{tab:calidad}, con el detalle completo en
\rc{logs/quality_report.csv}). Las reglas son:

\begin{itemize}
\item \textbf{Coordenada vac\'ia o en cero.} Antes de descartar, intento
recuperarla: si el registro trae un UBIGEO que corresponde a un distrito
real, ubico el establecimiento en el centro de ese distrito ---una regla
clara y reproducible sobre los pol\'igonos del INEI--- y lo marco como
ubicaci\'on aproximada. As\'i recuper\'e \coordRecuperadas{} de
\coordFaltantes{} registros sin coordenada (\textbf{tasa de recuperaci\'on
\tasaRecuperacion\,\%}). Entre los recuperados est\'an el Hospital II
Pucallpa (EsSalud) y el Hospital Intercultural de Atalaya ---este \'ultimo,
el \'unico establecimiento resolutivo de todo el sur de Ucayali.
\item \textbf{Signo del hemisferio cambiado.} Si una coordenada es positiva
pero su valor negativo cae dentro del Per\'u, le corrijo el signo. El Per\'u
est\'a en latitudes y longitudes negativas; una positiva es casi siempre un
error de tipeo.
\item \textbf{Latitud y longitud intercambiadas.} Si el par tal como viene
cae fuera del Per\'u pero al intercambiarlo cae dentro, lo intercambio.
\item \textbf{Fuera del Per\'u.} Lo que sigue fuera del recuadro del pa\'is
tras las correcciones anteriores se elimina, con la regla registrada.
\item \textbf{C\'odigos duplicados.} Me quedo con la primera aparici\'on.
\item \textbf{Punto fuera de su distrito.} No lo elimino ---puede ser el
pol\'igono el que est\'a mal, no el punto---; lo marco y lo vuelco a un
archivo aparte.
\item \textbf{Texto mal codificado} (``\~n'' que aparece como otro s\'imbolo).
Lo revierto.
\end{itemize}

\inputsihay{tables/calidad_datos.tex}

\subsection{C\'omo calcul\'e los tiempos de viaje}
Constru\'i un grafo de la red vial directamente desde el archivo local de
OpenStreetMap (\texttt{peru-latest.osm.pbf}) con la librer\'ia
\texttt{pyosmium}, y calcul\'e los caminos m\'as cortos con Dijkstra
vectorizado (\texttt{scipy.sparse.csgraph}).

\textbf{Decisi\'on 1: no us\'e OSRM en Docker}, que era la opci\'on
sugerida. La BIOS del equipo tiene la virtualizaci\'on deshabilitada y
bloqueada por el \'area de sistemas, as\'i que Docker y WSL2 no arrancan.
\textbf{Decisi\'on 2: descart\'e \texttt{pyrosm}} (otra opci\'on para leer
el \texttt{.pbf}) porque no compila en la versi\'on de Python del equipo.
\textbf{Decisi\'on 3: descart\'e OSMnx / Overpass} porque el servidor
p\'ublico estuvo inalcanzable los d\'ias de trabajo ---lo cual es, en s\'i,
un dato sobre el estado de las herramientas de datos abiertos. La soluci\'on
final ---leer el \texttt{.pbf} local--- no depende de internet ni de
permisos de administrador, y es completamente reproducible.

Par\'ametros: tres perfiles ---auto (velocidad seg\'un el tipo de v\'ia, y
la se\~nalizada cuando existe), bicicleta (15\,km/h) y a pie (4{,}5\,km/h).
Cada punto se ``engancha'' al nodo m\'as cercano de la red; si el nodo
m\'as cercano est\'a a m\'as de 1{,}5\,km, el punto se marca \textbf{sin
acceso por carretera} y no se rutea desde un nodo lejano. \textbf{En ning\'un
caso reemplazo un tramo faltante por una l\'inea recta.} La matriz completa
origen\,$\times$\,establecimiento (resolutivos m\'as los \mbox{I-3} y
\mbox{I-4}, que son los que el simulador de la Fase~4 permite ``ascender'')
se calcula para el perfil auto y se guarda en Parquet.

\subsection{Por qu\'e muestre\'e}
Los tres departamentos suman 17\,187 centros poblados y el enunciado pone un
tope de 5\,000. Tom\'e una muestra \textbf{estratificada por distrito}
(reparto proporcional, con una semilla fija para que sea reproducible). Como
SIGMED no trae poblaci\'on, la estratificaci\'on es por distrito y no por
poblaci\'on; la consecuencia (distritos con muchos caser\'ios peque\~nos
quedan algo sub-representados) la discuto en Limitaciones.

\subsection{Los indicadores}
Cada indicador es una funci\'on que recibe una tabla y devuelve una tabla
(\texttt{src/metrics.py}):
\begin{enumerate}
\item \textbf{Tiempo de acceso} $t_{\min}(i)$: minutos en auto desde el centro
poblado $i$ al resolutivo m\'as cercano.
\item \textbf{Bandas de cobertura}: qu\'e porcentaje de poblaci\'on est\'a a
menos de 30, entre 30 y 60, entre 60 y 120, a m\'as de 120 minutos, o sin
ruta.
\item \textbf{Acceso medio ponderado por poblaci\'on}, agregado por distrito,
provincia y departamento. El promedio se calcula solo sobre la poblaci\'on
que s\'i tiene ruta, y aparte se reporta cu\'anta poblaci\'on qued\'o sin
ruta.
\item \textbf{Lista de brechas cr\'iticas}: los distritos con peor acceso
ponderado.
\item \textbf{Desigualdad}: coeficiente de Gini y curva de Lorenz del tiempo
de acceso ponderado por poblaci\'on. Eleg\'i el Gini porque es un solo
n\'umero entre 0 y 1, f\'acil de comparar entre departamentos y contra
futuras mediciones, y es el est\'andar en estudios de equidad en el acceso.
\item \textbf{Urbano frente a rural}, con la regla del INEI: urbano si el
centro poblado tiene 2\,000 habitantes o m\'as, o es capital de distrito.
\end{enumerate}
Adem\'as, un cruce entre acceso y \textbf{altitud}, para ver si el tiempo de
viaje empeora con la altura.

\FloatBarrier
\section{Resultados}

De \poblacionTotal{} habitantes, el \pctCubierta\,\% llega a un establecimiento
resolutivo en \umbralOro{} minutos o menos y la mediana del tiempo de acceso
es de \medianaAcceso{} minutos; pero un \pctSinRuta\,\%
(\poblacionSinRuta{} personas), casi todo en Ucayali, vive sin ninguna
carretera hacia un hospital resolutivo. Lo que sigue desglosa ese panorama
por departamento (Cuadro~\ref{tab:acceso-dep}), por banda de tiempo
(Figura~\ref{fig:cobertura_bandas} y Cuadro~\ref{tab:cobertura}), urbano
frente a rural y por altitud (Figura~\ref{fig:ecdf_urbano_rural}), por
desigualdad (Figura~\ref{fig:lorenz}, Cuadro~\ref{tab:gini}) y por distrito
(Cuadros~\ref{tab:brechas} y~\ref{tab:urbano-rural}).

\inputsihay{tables/acceso_departamento.tex}

\figsihay[0.6\linewidth]{cobertura_bandas}{Porcentaje de poblaci\'on por
banda de tiempo de acceso en auto, por departamento. En Ucayali una parte
importante de la poblaci\'on cae en la banda ``sin ruta''.}

\inputsihay{tables/cobertura_bandas.tex}

\dosfiguras{ecdf_urbano_rural}{acceso_altitud}{\emph{Izquierda:} porcentaje
acumulado de poblaci\'on seg\'un el tiempo de acceso, urbano frente a rural;
la l\'inea punteada marca el umbral de \umbralOro{} minutos. \emph{Derecha:}
acceso medio ponderado por franja de altitud (correlaci\'on de Spearman entre
tiempo y altitud: \spearmanAltitud).}{}

\figsihay[0.34\linewidth]{lorenz}{Curva de Lorenz del tiempo de acceso
ponderado por poblaci\'on (Gini $=\giniTotal$). Cuanto m\'as se aleja de la
diagonal, m\'as desigual es el acceso.}

\inputsihay{tables/gini.tex}

Bajando al detalle, el Cuadro~\ref{tab:brechas} lista los 15 distritos con
peor acceso ponderado ---los cinco primeros, todos amaz\'onicos, con el
100\,\% de su poblaci\'on sin ruta--- y el Cuadro~\ref{tab:urbano-rural}
separa lo urbano de lo rural en cada departamento: la brecha urbano/rural es
sistem\'atica y m\'axima en Ucayali.

\inputsihay{tables/brechas_criticas.tex}

\inputsihay{tables/urbano_rural.tex}

\FloatBarrier
\section{Discusi\'on}

\subsection{L\'inea recta frente a carretera}
La Figura~\ref{fig:recta_vs_red} compara, para cada centro poblado con ruta,
la distancia en l\'inea recta a su resolutivo m\'as cercano con la distancia
real por carretera. El factor de rodeo mediano es \factorRodeo: la carretera
recorre, en la mitad de los casos, al menos \factorRodeo{} veces la l\'inea
recta. Y ese factor no es parejo ---es mayor en la sierra, donde la
carretera bordea el relieve. Por eso a los puntos sin ruta los dej\'e como
sin ruta: multiplicar una l\'inea recta por un factor fijo dar\'ia un
n\'umero que parece preciso pero no lo es.

\figsihay[0.46\linewidth]{recta_vs_red}{Distancia por carretera frente a
distancia en l\'inea recta al resolutivo m\'as cercano (perfil auto). La
recta marca el factor de rodeo mediano.}

\subsection{El contraste que importa}
Piura y Ucayali no se comportan igual, y no deber\'ian. En \textbf{Piura} la
red vial es densa y el acceso depende de la distancia a las pocas ciudades
con hospital \mbox{II-1}; el problema es de \emph{distancia}. En
\textbf{Ucayali} la mayor\'ia de la poblaci\'on vive en centros poblados a
la orilla de los r\'ios, sin carretera: el transporte real es fluvial, y un
an\'alisis vial ---como este--- lo refleja marcando a esa gente como sin
ruta por carretera; el problema es de \emph{modo} de transporte.
\textbf{Ayacucho} est\'a en el medio: hay v\'ias, pero el rodeo por el
relieve alarga mucho los tiempos. Un promedio distrital sin ponderar
esconder\'ia todo esto, porque tratar\'ia a un caser\'io de 12 personas
igual que a una ciudad de 12\,000; ponderar por poblaci\'on lo hace visible.

\FloatBarrier
\section{Limitaciones}

\begin{itemize}
\item \textbf{Errores del registro.} Del $\sim$37\,\% de RENIPRESS sin
coordenada recuper\'e \tasaRecuperacion\,\% al centro del distrito; el resto
se excluye. Si los excluidos no son al azar (por ejemplo, m\'as frecuentes en
zonas remotas), la cobertura puede estar sesgada.
\item \textbf{Las coordenadas recuperadas son aproximadas.} Los
establecimientos recuperados quedan en el centro de su distrito, no en su
ubicaci\'on real; en distritos amaz\'onicos grandes eso puede correr el punto
varios kil\'ometros. Van marcados y con borde distinto en el mapa.
\item \textbf{La red de OpenStreetMap es despareja.} En zonas rurales y
amaz\'onicas est\'a incompleta. Algunos puntos ``sin ruta'' quiz\'a tengan
una trocha no mapeada; otros, una v\'ia mapeada que en la pr\'actica no se
puede transitar. El resultado es una cota, no una medici\'on exacta.
\item \textbf{``Activo'' no es lo mismo que ``resolutivo en la pr\'actica''.}
Que un hospital figure como \mbox{II-1} activo no garantiza que esa noche
haya cirujano, anestesi\'ologo, banco de sangre y quir\'ofano disponibles.
\item \textbf{No hay ambulancias ni sistema de referencia en el modelo.}
Asumo que el paciente se traslada por sus propios medios.
\item \textbf{Muestreo.} Rute\'e cerca del \fraccionMuestra\,\% de los
centros poblados, estratificado por distrito y no por poblaci\'on. Las
cifras de cobertura son estimaciones con error de muestreo.
\item \textbf{La poblaci\'on es modelada.} WorldPop es un modelo a 1\,km, no
un censo; ajust\'e los totales al pol\'igono departamental para que cuadren
con el censo, pero la distribuci\'on fina dentro del distrito es aproximada.
\item \textbf{Velocidades supuestas.} Las velocidades por tipo de v\'ia son
valores de referencia; no consideran el estado del asfalto, la \'epoca de
lluvias ni el tr\'afico.
\end{itemize}

\FloatBarrier
\section{Conclusiones y recomendaciones}

\begin{enumerate}
\item El acceso a atenci\'on resolutiva en los tres departamentos es muy
desigual (Gini $=\giniTotal$), y esa desigualdad es \emph{geogr\'afica}. La
Amazon\'ia no necesita m\'as carreteras: necesita capacidad resolutiva por
r\'io ---lanchas-ambulancia, telemedicina en los puestos ribere\~nos--- y
por eso la recomendaci\'on para Ucayali va por ah\'i.
\item En Ayacucho, donde el problema es el rodeo por el relieve, s\'i sirven
las mejoras de v\'ia (asfaltado, puentes). El simulador de la Fase~4 permite
probar qu\'e establecimiento \mbox{I-3} o \mbox{I-4} conviene ascender para
cubrir m\'as poblaci\'on con la menor inversi\'on.
\item El distrito con peor acceso ponderado (\peorDistrito) deber\'ia
encabezar cualquier plan de inversi\'on.
\item Antes de decidir nada, hay que sanear RENIPRESS: el 37\,\% de
registros sin coordenada es la mayor fuente de incertidumbre de todo el
an\'alisis.
\end{enumerate}

\section*{Referencias}
\begingroup\small
\begin{itemize}\setlength{\itemsep}{1pt}
\item SUSALUD. Registro Nacional de IPRESS (RENIPRESS). \url{https://www.datosabiertos.gob.pe/dataset/registro-nacional-de-entidades-prestadoras-de-servicios-de-salud-renipress}
\item MINEDU. Descarga espacial SIGMED. \url{https://sigmed.minedu.gob.pe/descargas/}
\item OpenStreetMap contributors / Geofabrik. Extracto de Per\'u. \url{https://download.geofabrik.de/south-america/peru-latest.osm.pbf}
\item INEI. Cartograf\'ia censal (l\'imites pol\'itico-administrativos).
\item WorldPop (2020). Global High Resolution Population Denominators Project, Per\'u, ajustado a totales ONU. \rc{DOI:10.5258/SOTON/WP00685}.
\item Hagberg, A. \emph{et al.} NetworkX. \url{https://networkx.org}
\end{itemize}
\endgroup

\end{document}
